\section{Attribution Pipeline}
\label{sec:attribution}

\sys ingests two data sources produced for the same workload run: (1) an \nsys SQLite export
containing CUDA kernel activity and NVTX range annotations, and (2) \ncu CSV output containing
per-kernel hardware counter measurements. The pipeline produces a hardware-attributed
\texttt{OperatorProfile}: a JSON document mapping each \texttt{aten::} operator to its constituent
kernels, their hardware metrics, and the confidence tier of each attribution.

\subsection{Data Ingestion}
\label{sec:ingestion}

\textbf{Nsight Systems export.}  Running the workload under
\texttt{nsys profile --trace=cuda,nvtx} produces a \texttt{.nsys-rep} file, exported to SQLite
via \texttt{nsys export --type=sqlite}. \sys queries three tables:

\begin{lstlisting}[style=sqlstyle,caption={Kernel and NVTX range extraction from the \nsys SQLite export.}]
-- CUDA kernel activity (GPU-domain timestamps)
SELECT k.correlationId, s.value AS kernel_name,
       k.start AS start_ns, k.end AS end_ns,
       k.streamId, k.deviceId,
       k.gridX, k.gridY, k.gridZ,
       k.blockX, k.blockY, k.blockZ,
       r.globalTid AS host_tid,
       r.start      AS cpu_launch_start_ns
FROM   CUPTI_ACTIVITY_KIND_KERNEL k
JOIN   StringIds s ON s.id = k.shortName
LEFT JOIN CUPTI_ACTIVITY_KIND_RUNTIME r
       ON r.correlationId = k.correlationId
ORDER  BY k.start ASC;

-- NVTX range annotations (CPU-domain timestamps)
SELECT n.text, n.start AS start_ns, n.end AS end_ns,
       n.nestingLevel, n.globalTid AS stream_id
FROM   NVTX_EVENTS n
WHERE  n.end IS NOT NULL
ORDER  BY n.start ASC;
\end{lstlisting}

The \texttt{cpu\_launch\_start\_ns} field from the RUNTIME table is the host-side timestamp when
the CUDA kernel launch API returned, providing a CPU-domain anchor for NVTX bracket queries
without conflating it with the GPU-domain kernel execution timestamp (\secref{sec:clockdomain}).

\textbf{Nsight Compute export.}  For the full set of kernels in the \nsys trace, \sys invokes
a single application-mode replay pass per counter group (4--8 passes total), collecting all
20 metrics simultaneously:
\begin{lstlisting}[style=pystyle]
ncu --replay-mode application \
    --metrics <comma_separated_list> \
    --export <output.ncu-rep> \
    [--kernel-name <filter>] \
    <replay_script.py>
\end{lstlisting}
Application-mode replays the workload once per counter group regardless of the number of
unique kernel names, making profiling cost $O(\text{counter groups})$ rather than
$O(\text{unique kernels})$---a 10--50$\times$ efficiency gain for multi-layer models.
The optional \texttt{--kernel-name} flag filters to a specific kernel for targeted
re-profiling.  The resulting \texttt{.ncu-rep} is exported to CSV
(\texttt{ncu --import --csv}), producing one row per kernel invocation per metric.

\subsection{GPU Clock Locking for Reproducible Duration Comparisons}
\label{sec:clocklocking}

Per-operator durations in \texttt{profile.json} are derived from \nsys CUPTI GPU-domain
timestamps.  These durations are subject to a subtle confound: a GPU's boost clock floats
with die temperature and power draw, typically varying 5--15\% across captures separated
by minutes or hours.  An optimized workload dissipates less power than the baseline (fewer
FLOPs on the SIMT path), so its capture may run at a \emph{higher} sustained clock,
inflating the apparent speedup.  Conversely, thermal throttling during a long baseline
capture deflates the baseline.  Both effects corrupt the $\text{baseline}/\text{optimized}$
duration ratio even when no hardware change was made.

\textbf{Probe-and-lock.}
\sys pins the SM and memory clocks for the duration of the \nsys capture.  The default
target is \emph{probe-and-lock}: a synthetic matrix-multiply load runs for ${\sim}1.5\,\text{s}$
(with a $1\,\text{s}$ warmup to reach thermal steady state), and the median clock observed
during the probe is selected as the lock target.  Locking to the \emph{sustained} frequency
(rather than the peak boost clock, which the GPU cannot maintain under load) guarantees the
lock is genuinely held throughout the capture.

\textbf{Cache and reuse.}
The probed clock is written to \texttt{profiler\_output/.gpu\_clock\_lock.json} with a
6-hour TTL per GPU index.  The baseline capture probes once; the optimized capture reuses
the cached clock.  Both captures therefore execute at exactly the same frequency, making
their duration ratio clock-immune by construction---the same guarantee that invocation-order
matching provides for the hardware counter join.

\textbf{Implementation.}
Clocks are set via \texttt{nvidia-smi -lgc/-lmc} and released through a context manager
that registers a \texttt{SIGTERM} handler, guaranteeing a reset even on exception, timeout,
or Ctrl-C:
\begin{lstlisting}[style=pystyle]
with gpu_clocks_locked(target="probe", gpu_index=0,
                       use_sudo=True, cache_dir=cache_dir):
    # nsys profile ... run_workload.py runs here
    ...
\end{lstlisting}
Clock setting uses \texttt{sudo -n} (non-interactive); if permission is unavailable, \sys
logs a \texttt{WARNING} and proceeds at dynamic clocks rather than aborting.

\textbf{Scope.}
Clock locking applies to the \nsys capture phase only.  The \ncu replay phase is
unaffected: \ncu already self-locks to its own base clock before each replay pass, and
the hardware counter values it reports are independent of the GPU's runtime SM frequency.

\subsection{The Clock Domain Problem and Invocation-Order Matching}
\label{sec:clockdomain}

The fundamental obstacle in cross-tool profiling is that \nsys and \ncu observe the same kernel
execution through incompatible timing mechanisms. \nsys records CUPTI GPU-domain timestamps from
the hardware's free-running clock. \ncu replays kernels in an isolated profiling pass with injected
performance counter reads; its internal invocation clock reflects replay-mode execution, not the
original capture clock. On a single GPU, these two clocks routinely diverge by 10--100~$\mu$s per
kernel due to replay serialization and PMU injection overhead. A naive timestamp join would
incorrectly associate hardware counter measurements with the wrong operator---especially dangerous
for operators with similar execution times.

\textbf{Invocation-Order Matching (Algorithm~\ref{alg:iom}).}  The key insight is that \ncu
replays kernels in the same execution order they appear in the \nsys timeline: the $i$-th
invocation of kernel $K$ in the \nsys trace corresponds to the $i$-th row for kernel $K$ in the
\ncu CSV output.  We exploit this ordering to assign hardware metrics without any timestamp
comparison:

\begin{algorithm}[t]
\caption{Invocation-Order Matching (\nsys $\times$ \ncu)}
\label{alg:iom}
\begin{algorithmic}[1]
\Require \texttt{nsys\_kernels}: list of \texttt{KernelRow} sorted by \texttt{start\_ns} per stream
\Require \texttt{ncu\_rows}: list of \texttt{NcuRow} sorted by invocation index per kernel name
\Ensure  Each \texttt{KernelRow} assigned a \texttt{KernelMetrics} object
\State $\textit{counter} \gets \text{defaultdict}(\text{int})$
  \Comment{per kernel name}
\For{$k \in \texttt{nsys\_kernels}$}
  \State $i \gets \textit{counter}[k.\textit{name}]$
  \If{$i \ge |\texttt{ncu\_rows}[k.\textit{name}]|$}
    \State \textbf{raise} \texttt{InvocationCountMismatchError}($k.\textit{name}$, $i$)
    \Comment{iteration count mismatch between captures}
  \EndIf
  \State $k.\textit{metrics} \gets \texttt{ncu\_rows}[k.\textit{name}][i]$
  \State $\textit{counter}[k.\textit{name}] \mathrel{+}= 1$
\EndFor
\end{algorithmic}
\end{algorithm}

This matching is correct when the workload is \emph{deterministic}: the same kernel names must
appear in the same counts and the same execution order in both the \nsys capture and the \ncu
replay.  Concretely, models with conditional control flow, dynamic dispatch, or
\texttt{cudaMallocAsync} workspace kernels can violate this invariant without producing a shape
error.  \sys validates that per-kernel invocation counts match between the two runs before
issuing counter assignments; a count mismatch on any kernel name raises
\texttt{InvocationCountMismatchError} and aborts attribution for that workload.  Shape validation
before the replay step enforces the shape portion of this invariant
(\secref{sec:edgecases}).  \figref{fig:iom} illustrates why timestamp joins fail and how
invocation-order matching avoids the problem.

\begin{figure}[t]
\centering
\begin{tikzpicture}[
  font=\small,
  arr/.style={-{Stealth[length=4pt]}, thick},
  karr/.style={-{Stealth[length=4pt]}, thick, red!70!black},
  garr/.style={-{Stealth[length=4pt]}, thick, green!50!black},
  kern/.style={draw, rounded corners=2pt, fill=blue!10, draw=blue!60!black,
               minimum width=1.6cm, minimum height=1.4em, align=center, font=\footnotesize},
  kernncu/.style={draw, rounded corners=2pt, fill=purple!10, draw=purple!60!black,
               minimum width=1.6cm, minimum height=1.4em, align=center, font=\footnotesize},
  tl/.style={-{Stealth[length=5pt]}, thick, gray},
  lbl/.style={font=\scriptsize, text=gray!80!black},
  hl/.style={font=\footnotesize\bfseries},
  darr/.style={-{Stealth[length=4pt]}, thick, dashed, gray!50},
]

%% ---- (a) Naive timestamp join ----
\node[hl] at (4.5, 0.7) {(a) Naive timestamp join --- \textcolor{red!70!black}{incorrect}};

\draw[tl] (0.0, 0) -- (9.8, 0) node[right,lbl] {nsys (GPU clock)};
\node[lbl, left=2pt] at (0.0, 0) {\small nsys};

\node[kern] (k1n) at (1.1, 0)  {$K_1$\;\tiny t=10};
\node[kern] (k2n) at (4.0, 0)  {$K_2$\;\tiny t=45};
\node[kern] (k3n) at (7.1, 0)  {$K_3$\;\tiny t=80};

\draw[tl] (0.0,-1.25) -- (9.8,-1.25) node[right,lbl] {ncu (replay clock)};
\node[lbl, left=2pt] at (0.0,-1.25) {\small ncu};

\node[kernncu] (k1c) at (2.3,-1.25)  {$K_1$\;\tiny t=35};
\node[kernncu] (k2c) at (5.1,-1.25)  {$K_2$\;\tiny t=62};
\node[kernncu] (k3c) at (8.1,-1.25)  {$K_3$\;\tiny t=95};

\draw[darr] (k1c.north) -- (k1n.south);
\draw[karr] (k2c.north) -- (k1n.south)
  node[near start, left=3pt, lbl, text=red!70!black] {\textbf{\ding{55}} wrong match};
\draw[darr] (k3c.north) -- (k3n.south);

\draw[decorate,decoration={brace,amplitude=5pt},gray!60]
  (k1n.south east)++(0.1,0) -- (k1c.north east)++(0.1,0)
  node[midway, right=6pt, lbl, align=left] {clock\\drift};

%% ---- (b) IOM ----
\node[hl] at (4.5,-2.15) {(b) Invocation-order matching (IOM) --- \textcolor{green!50!black}{correct}};

\draw[tl] (0.0,-2.9) -- (9.8,-2.9) node[right,lbl] {nsys order};
\node[lbl, left=2pt] at (0.0,-2.9) {\small nsys};
\node[kern] (k1ni) at (1.4,-2.9) {$K_1$\;\tiny $i{=}0$};
\node[kern] (k2ni) at (4.0,-2.9) {$K_2$\;\tiny $i{=}1$};
\node[kern] (k3ni) at (6.7,-2.9) {$K_3$\;\tiny $i{=}2$};

\draw[tl] (0.0,-4.1) -- (9.8,-4.1) node[right,lbl] {ncu order};
\node[lbl, left=2pt] at (0.0,-4.1) {\small ncu};
\node[kernncu] (k1ci) at (1.4,-4.1) {$K_1$\;\tiny $i{=}0$};
\node[kernncu] (k2ci) at (4.0,-4.1) {$K_2$\;\tiny $i{=}1$};
\node[kernncu] (k3ci) at (6.7,-4.1) {$K_3$\;\tiny $i{=}2$};

\draw[garr] (k1ci.north) -- (k1ni.south) node[midway,right=1pt,lbl,text=green!50!black] {\textbf{\ding{51}}};
\draw[garr] (k2ci.north) -- (k2ni.south) node[midway,right=1pt,lbl,text=green!50!black] {\textbf{\ding{51}}};
\draw[garr] (k3ci.north) -- (k3ni.south) node[midway,right=1pt,lbl,text=green!50!black] {\textbf{\ding{51}}};

\node[lbl, align=center] at (4.9,-4.75)
  {Match key: \texttt{(kernel\_name, invocation\_index)} --- no timestamp comparison required};

\end{tikzpicture}
\caption{\textbf{Why timestamp joins fail and how IOM avoids them.}
  (a)~The nsys GPU-domain clock and the \ncu replay clock diverge by $>$10\,$\mu$s; a
  nearest-timestamp join misattributes $K_2$'s hardware counters to $K_1$.
  (b)~IOM matches kernels by ordinal position within each \texttt{(kernel\_name,
  invocation\_index)} sequence, making attribution clock-immune by construction.}
\label{fig:iom}
\end{figure}

\subsection{StreamIntervalTree Attribution}
\label{sec:intervaltree}

For each unique \texttt{(stream\_id, device\_id)} pair observed in the \nsys trace, \sys
constructs a \texttt{StreamIntervalTree}: a sorted list of NVTX ranges indexed by
\texttt{start\_ns}, implemented using Python's \texttt{sortedcontainers.SortedList} for
$O(\log n)$ insertion and $O(k \log n)$ enclosure queries ($k$ = number of enclosing ranges).

\begin{algorithm}[t]
\caption{NVTX Enclosure Attribution (MEDIUM Confidence)}
\label{alg:nvtx}
\begin{algorithmic}[1]
\Require Kernel row $k$ with \texttt{cpu\_launch\_start\_ns}, \texttt{stream\_id}
\Require \texttt{forest}: map from \texttt{(stream\_id, device\_id)} to \texttt{SortedList}
\Ensure Attribution at MEDIUM confidence, or fall through to heuristic
\State $T \gets k.\textit{cpu\_launch\_start\_ns}$
\State $\textit{ranges} \gets$ \texttt{forest}[$k.\textit{stream\_id}$].query\_enclosing($T$)
  \Comment{all NVTX ranges $[s,e]$ with $s \le T \le e$}
\If{$|\textit{ranges}| = 0$}
  \State \textbf{fall through} to Inductor fusion enrichment (\secref{sec:heuristic}) or UNATTRIBUTED
\ElsIf{$|\textit{ranges}| = 1$}
  \State $k.\textit{operator} \gets \textit{ranges}[0].\textit{text}$,
         \textit{confidence} $\gets$ MEDIUM
\Else
  \State $k.\textit{operator} \gets \arg\max_r r.\textit{nesting\_level}$
         \Comment{innermost enclosing range; ties broken by latest \texttt{start\_ns}}
  \State $k.\textit{is\_fused} \gets \top$;
         $k.\textit{all\_enclosing} \gets \textit{ranges}$
         \Comment{fused kernel spans multiple operators}
\EndIf
\end{algorithmic}
\end{algorithm}

The \texttt{cpu\_launch\_start\_ns} field is used as the query point rather than the GPU
\texttt{start\_ns}, because NVTX ranges are CPU-side events and the async gap between host launch
and GPU execution makes GPU timestamps unreliable for enclosure queries.

\subsection{Inductor Fusion Enrichment}
\label{sec:heuristic}

When NVTX enclosure fails (no enclosing range exists at the query point), \sys applies a
post-attribution enrichment pass that recovers ground-truth fused operator identity from Inductor
debug artifacts, rather than guessing from kernel names.

\textbf{Inductor debug extraction.}  When \texttt{torch.compile(backend="inductor")} runs with
debug mode enabled, Inductor writes \texttt{output\_code.py} artifacts containing the generated
Triton kernel source.  Each kernel call site is preceded by a structured comment:

\begin{lstlisting}[style=pystyle]
# Topologically Sorted Source Nodes: [...],
# Original ATen: [aten.relu, aten.addmm]
def triton_poi_fused_relu_addmm_0(in_ptr0, ...):
    ...
\end{lstlisting}

\sys's \texttt{InductorFusionExtractor} parses these comments via the regex
\texttt{\# .+Original~ATen:\textbackslash{}s*\textbackslash{}[([\textasciicircum{}\textbackslash{}]]+)\textbackslash{}]},
producing a mapping $\{$\textit{kernel\_short\_name} $\to$ $[$\texttt{aten::op}_1,
\texttt{aten::op}_2, $\ldots]$$\}$ that contains the exact set of aten operators fused into
each kernel.

\textbf{Enrichment as a post-attribution pass.}  This map is applied \emph{after} the NVTX
enclosure pass and never overrides a valid TORCH\_PROFILER or NVTX attribution:
\begin{enumerate}[leftmargin=1.5em]
  \item For kernels that remain UNATTRIBUTED after NVTX enclosure, a lookup in the fusion map
    upgrades them to \texttt{INDUCTOR\_FUSION} attribution at MEDIUM confidence, with
    \texttt{source\_operators} populated from the \texttt{Original~ATen} comment.
  \item For already-attributed kernels (TORCH\_PROFILER or NVTX), the enrichment augments the
    \texttt{fused\_ops} field and sets \texttt{is\_fused~=~True} if multiple operators are present.
\end{enumerate}

Kernels matching neither NVTX enclosure nor the Inductor fusion map remain in
\texttt{unattributed\_kernels[]} with \texttt{attribution\_method = "unattributed"}.  For eager-mode
workloads or cuDNN-backed operators (\eg \texttt{nn.LSTM} via \texttt{aten::\_cudnn\_rnn}), which
produce neither NVTX brackets nor Inductor debug output, unattributed rates of 20--85\% are expected
and are surfaced explicitly in the profile output.

\subsection{Confidence Tiers}
\label{sec:tiers}

\begin{table}[t]
\centering
\caption{Attribution methods in \sys.  Coverage percentages are measured over kernel
  \emph{runtime} (not kernel count); the UNATTRIBUTED rate varies widely with compile mode and
  operator backend.  Ranges shown are typical across PyTorch-native workloads; see
  \tabref{tab:coverage} for per-workload values.}
\label{tab:tiers}
\small
\setlength{\tabcolsep}{4pt}
\begin{tabular}{p{2.0cm}p{2.6cm}p{4.8cm}p{1.6cm}}
\toprule
\textbf{Method} & \textbf{Evidence} & \textbf{Trigger Condition} & \textbf{Coverage} \\
\midrule
\rowcolor{highcol!15}
HIGH & torch.profiler IOM & Kineto trace available; kernel count agrees with dispatch count & 35--45\% \\
\rowcolor{medcol!15}
MEDIUM (NVTX) & NVTX enclosure (interval tree) & Kernel launch falls within an \texttt{aten::} NVTX range & 45--55\% \\
\rowcolor{medcol!10}
MEDIUM (Inductor) & Inductor \texttt{output\_code.py} comments & Kernel name in Inductor fusion map; NVTX enclosure absent & 5--15\% \\
\rowcolor{ucol!15}
UNATTRIBUTED & None & cuDNN-internal kernels (LSTM, cuDNN conv.); eager Triton without NVTX & 1--85\% \\
\bottomrule
\end{tabular}
\end{table}

\subsection{Edge Cases}
\label{sec:edgecases}

\sys explicitly handles the following attribution edge cases:

\begin{enumerate}[leftmargin=1.5em]

\item \textbf{Clock domain mismatch.} Invocation-order matching (\algref{alg:iom}) makes
  attribution clock-immune by construction.

\item \textbf{CUDA Graph replay.} During graph capture, operators emit NVTX ranges and kernels are
  attributed normally.  During replay, new GPU timestamps with no NVTX correspondence are generated.
  \sys detects replay by the absence of NVTX events in the replay window and replays the attribution
  cache from the capture phase.

\item \textbf{Multi-stream execution.} Each \texttt{(stream\_id, device\_id)} pair has its own
  \texttt{StreamIntervalTree}, preventing cross-stream false matches.

\item \textbf{JIT warm-up kernels.} Kernels launched before the first NVTX range (JIT compilation
  artifacts) are flagged using a duration-outlier heuristic ($> 10\times$ median kernel duration)
  and excluded from operator mapping.

\item \textbf{Async kernel launch.} NVTX push/pop occurs on the host before the kernel executes on
  the GPU.  The RUNTIME table's \texttt{cpu\_launch\_start\_ns} field provides a host-side anchor
  that lies within the NVTX range, whereas the GPU \texttt{start\_ns} may not.

\item \textbf{Dynamic shapes.} \sys validates that input shapes at \ncu replay match those recorded
  during \nsys capture; a \texttt{ShapeMismatchError} is raised otherwise.

\item \textbf{Fused kernels.} When multiple NVTX ranges enclose a single kernel (Triton-fused
  cross-operator), \sys sets \texttt{is\_fused=True} and stores all enclosing ranges in
  \texttt{all\_enclosing\_ranges} for downstream analysis.

\item \textbf{ncu timing independence.} \ncu timestamps are internal to each replay run and are
  discarded entirely; only metric values are extracted from \ncu output.

\end{enumerate}

\subsection{Layer Deduplication}
\label{sec:dedup}

For models with repeated structure (\eg transformer decoder blocks), profiling and
compiling every layer independently costs $O(N)$ ncu replay time.
\texttt{UniqueSubgraphRegistry} detects structurally identical FX subgraphs by
computing a canonical hash of each partition's node types and connectivity.  Only
one representative per equivalence class is compiled and profiled; metrics and
compiled callables are propagated to all structural duplicates by positional index
within each partition group.

The deduplication is transparent to the outer profiling loop: each duplicate
partition's forward pass is wrapped with an NVTX range tagged
\texttt{layer::duplicate::<label>}, and the partition equivalence map (written to
\texttt{<prefix>.part.json} during the correlation pass) is passed to
\texttt{KernelProfileConfig} so the \ncu orchestrator skips duplicate-partition
kernels during replay and propagates metrics from the unique representative.

\textbf{Impact.}  For GPT-2 (12 identical transformer blocks),
deduplication reduces \ncu replay time from ${\sim}$30--45\,min to
${\sim}$3--5\,min---a 12$\times$ improvement.  Layer deduplication is activated
automatically by \texttt{run\_workload.py} when structurally repeated partitions
are detected.

\textbf{Warmup-gating fix.}  An important correctness invariant: the
\texttt{layer::unique::} and \texttt{layer::duplicate::} NVTX ranges must only be
emitted during the \emph{measured} capture window, not during warmup iterations.
\texttt{run\_workload.py} gates all layer-partition NVTX emission behind an
\texttt{\_NVTX\_STATE[`active']} flag that is set only after warmup completes.
Without this guard, warmup-phase kernels appear in the nsys trace inside
layer-partition ranges, causing \texttt{ManifestBuilder} to misattribute them to
operators rather than classifying them as pre-NVTX initialization---inflating
attributed kernel counts and biasing operator durations.

% -------------------------------------------------------
\section{Metric Set and Aggregation}
\label{sec:metrics}

\subsection{Counter Selection Rationale}

Collecting all 90+ \ncu counters is expensive (multiple replay passes) and produces profiles too
large for operator-level analysis.  We select 20 counters according to three criteria:
(1)~coverage of every bottleneck axis in the roofline taxonomy (\secref{sec:roofline});
(2)~availability of a valid counter name on Ampere, Hopper, \emph{and} Blackwell (some counters
were renamed in the Blackwell GB202 silicon); (3)~non-redundancy (no two selected counters
correlate perfectly).

\subsection{The 20-Counter Set}

\sys collects 20 hardware counters covering all major bottleneck axes: memory bandwidth,
compute throughput, Tensor Core utilization, occupancy, latency, and register pressure.
The full counter list with bottleneck axis, aggregation method, and architecture-specific
name fallbacks is in \tabref{tab:counters} (\secref{sec:appendix:counters}).  Seven counters
require name fallbacks between Ampere, Hopper, and Blackwell; \sys queries all variants and
uses the first that returns a value.

\subsection{Duration-Weighted Aggregation}
\label{sec:aggregation}

An operator may dispatch multiple kernels with heterogeneous durations.  For rate and utilization
metrics (rows marked DW-mean in \tabref{tab:counters}), a naive arithmetic mean would give a
short-lived auxiliary kernel the same weight as the long-running compute kernel.  \sys uses
duration-weighted aggregation instead:
\begin{equation}
  \overline{m}_\text{op} = \frac{\displaystyle\sum_{i} m_i \cdot d_i}{\displaystyle\sum_{i} d_i}
  \label{eq:dwmean}
\end{equation}
where $m_i$ and $d_i$ are the metric value and duration of the $i$-th kernel attributed to the
operator.  This ensures that the dominant kernel---the primary optimization target---contributes
proportionally to the operator-level summary.  For additive quantities (bytes, instruction counts,
spills), \sys uses plain summation.  For per-kernel constants (register count, shared memory),
\sys takes the maximum across all kernels, since the highest-pressure kernel limits occupancy.

% -------------------------------------------------------
\section{Applying the Profile: Optimization Case Study}
\label{sec:pgo}

To demonstrate that the attributed hardware profiles produced in \secref{sec:attribution}
are actionable, we implement a profile-guided optimization workflow in which per-operator
counter values from \texttt{profile.json} drive the selection of FX graph transformations.
This section describes the optimization infrastructure and patterns; their empirical results
are in \secref{sec:experiments}.  The optimization layer is one possible downstream
application of the attribution pipeline---others include hardware-aware neural architecture
search, performance regression detection in CI, and kernel-level profiling of serving stacks
such as vLLM or SGLang.

\subsection{FX Pass Infrastructure}

\sys provides two composable components for graph-level optimization.
\textbf{UniqueSubgraphRegistry} partitions the FX graph by detected layer
structure, identifies structurally identical subgraphs via hash comparison,
and compiles a single representative per equivalence class.  Duplicate
partitions share the compiled callable, reducing compilation and ncu replay
time by $O(\text{num\_duplicate\_layers})$.

\textbf{FxPassRunner} applies user-defined pattern rewrites to unique
representatives only, then propagates the transformed compiled callables to
all structural duplicates:

\begin{lstlisting}[style=pystyle]
runner = FxPassRunner(registry)
for pattern_fn, replacement_fn in get_passes():
    runner.apply_pass(pattern_fn, replacement_fn)
\end{lstlisting}

Custom backends are registered via Torch Dynamo's backend hook and receive
the Aten IR \texttt{GraphModule} before Inductor lowering.

\subsection{Three-Level FX Pass Architecture}
\label{sec:threelevel}

Profile-guided optimizations are routed across three IR levels depending on the
graph representation at which each transformation is expressible.

\textbf{Level 1 --- Functional (pre-AOTAutograd).}
Fusion and operator substitution passes run here.  The key invariant is that the
Dynamo graph preserves the single shared-activation node for operations like a QKV
projection: at this level, three linear layers reading from one \texttt{LayerNorm}
output form a three-GEMM pattern that can be fused into one wide GEMM.
AOTAutograd decomposition shatters that shared node into per-consumer
\texttt{aten.view} copies, making the pattern undetectable at the Aten level.
SDPA canonicalization similarly requires the high-level \texttt{F.scaled\_dot\_product\_attention}
call node, which is replaced by Inductor-specific primitives after decomposition.

\textbf{Level 2 --- Aten IR (\texttt{\_aten\_inner\_compile}).}
Dtype casts, BatchNorm folding, and channels-last annotation run against the fully
decomposed \texttt{torch.ops.aten.*} node set.  An important subtlety: the
\texttt{example\_inputs} argument of \texttt{inner\_compile} may be \texttt{FakeTensor}
objects (symbolic values without storage), making it impossible to read weight values
directly.  Weight-reading passes (BN fold, pre-transposition) receive \texttt{real\_inputs}
threaded through \texttt{functools.partial} at the outer backend call site, before any
compilation begins.

\textbf{Level 3 --- Inductor config (\texttt{config\_patches}).}
Non-graph directives---\texttt{freezing}, \texttt{max\_autotune}, \texttt{max\_autotune\_gemm}---
are passed as a scoped configuration dictionary to \texttt{compile\_fx} rather than
mutating the global \texttt{torch.\_inductor.config}.  This ensures the policy applies
only to the specific compilation and cannot leak across compilation units.

The three levels compose as a fixed funnel:
\begin{lstlisting}[style=pystyle]
@register_backend
def my_backend(gm, example_inputs):
    gm = _run_functional_passes(gm)               # Level 1
    inner = functools.partial(_aten_inner_compile,
                              real_inputs=list(example_inputs))
    return compile_fx(gm, example_inputs,          # Level 2 + 3
                      inner_compile=inner,
                      config_patches={"freezing": True})
\end{lstlisting}

Each pass degrades gracefully: a \texttt{try}/\texttt{except} wrapper logs
\texttt{WARNING} and returns the graph unchanged if the target pattern is absent,
so partial optimization is always preferred over a compilation failure.
\texttt{gm.graph.lint()} is called after each mutation to validate graph invariants.
When \texttt{UniqueSubgraphRegistry} detects repeated layer structure, the funnel
runs once per unique representative and shares the compiled callable across all
structural duplicates---each FX pass is therefore authored once for the full model.

\subsection{Optimization Patterns}
\label{sec:passes}

Hardware-evidence-motivated optimization patterns are implemented per-workload
in each example's \texttt{*\_optimized.py}.  Patterns fall into two categories:
\emph{pre-compile model modifications} applied before \texttt{torch.compile} is
called, and \emph{FX graph pattern rewrites} applied inside a custom backend via
\texttt{torch.fx.subgraph\_rewriter.replace\_pattern}.  Each pattern fires only
when the profiler has diagnosed the specific counter symptom that motivates it.
\tabref{tab:passes} summarizes the patterns demonstrated across the evaluated
workloads.

\begin{table}[t]
\centering
\caption{Hardware-evidence-motivated optimization patterns demonstrated across
  the evaluated workloads.  ``Pre-compile'' patterns modify the \texttt{nn.Module}
  or tensor before \texttt{torch.compile}; ``FX rewrite'' patterns modify the
  Aten IR graph inside a custom backend.}
\label{tab:passes}
\small
\setlength{\tabcolsep}{4pt}
\begin{tabular}{p{2.0cm}p{1.5cm}p{2.8cm}p{3.3cm}}
\toprule
\textbf{Pattern} & \textbf{Type} & \textbf{Counter Evidence} & \textbf{Transformation} \\
\midrule
BF16 promotion
  & Pre-compile
  & \texttt{tensor\_core\_active\_pct < 5\%}
  & \texttt{model.to(bfloat16)} before compile; routes GEMM to HMMA Tensor Core path \\[4pt]
\texttt{channels\_last} layout
  & Pre-compile
  & \texttt{convertTensor\_kernel} in \nsys trace
  & \texttt{model.to(channels\_last)}; eliminates cuDNN NCHW$\to$NHWC coercion \\[4pt]
BatchNorm folding
  & Pre-compile
  & \texttt{\_native\_batch\_norm\_legit\_no\_training} in profile
  & Fold BN running statistics into preceding Conv2d weights; replace module with \texttt{Identity} \\[4pt]
Weight pretranspose
  & FX rewrite
  & Low \texttt{sm\_throughput\_pct} on TN GEMM
  & Replace \texttt{aten.t(W)} + \texttt{aten.mm} with pre-computed \texttt{W$^\top$.contiguous()} buffer \\[4pt]
max-autotune
  & Compiler setting
  & Sub-optimal tile selection (SM throughput below peak)
  & \texttt{config\_patches=\{`max\_autotune': True\}} in \texttt{compile\_fx} \\
\bottomrule
\end{tabular}
\end{table}


\subsection{Closed-Loop Validation Protocol}
\label{sec:validation}

The full optimization cycle consists of four steps:
\begin{enumerate}
  \item \textbf{Profile}: capture baseline \texttt{profile.json} via \nsys + \ncu pipeline. The \nsys capture runs with GPU clocks locked to the probe-and-cache frequency (\secref{sec:clocklocking}); the same cached clock is reused for the optimized re-capture in step~4, ensuring duration ratios reflect architectural changes rather than thermal variation.
  \item \textbf{Diagnose}: \texttt{profile.json} is examined per-operator; per-operator hardware
    counter values are used to produce ranked optimization proposals (\texttt{optimizations.json})
    with evidence citations and dependency ordering.
  \item \textbf{Transform}: implement a workload-specific optimized backend and
    run \texttt{torch.compile(model, backend="<workload>\_opt")} to apply the patterns and compile.
  \item \textbf{Validate}: capture optimized \texttt{profile\_optimized.json} and generate a delta
    report comparing per-operator counter values, flagging any regression.
\end{enumerate}

This protocol makes every optimization auditable: the delta report shows not just speedup, but
\emph{which counter improved}, providing evidence that the diagnosed bottleneck was actually
eliminated rather than merely shifted to a different operator.

% -------------------------------------------------------
\section{Design Challenges}
\label{sec:challenges}

Implementing \sys required resolving several non-obvious engineering obstacles.
We document them here as they motivate the design decisions in \secref{sec:attribution}
and \secref{sec:pgo}, and because the failure modes are not documented elsewhere.

\subsection*{CUPTI Exclusivity: Two-Phase Capture}

\nsys and \texttt{torch.profiler} both consume CUPTI for kernel timing and cannot run simultaneously within the same process.  An early design ran the torch.profiler correlation pass \emph{inside} an \texttt{nsys profile} invocation, expecting both to collect independently.  In practice, CUPTI preempts the \texttt{torch.profiler} callback with no error; the correlation pass produces a \texttt{.corr.json} with zero entries, and the silence makes the failure invisible.  The pipeline was redesigned as two sequential invocations: Phase~1 is a plain Python run (no nsys) that captures torch.profiler output and writes \texttt{.corr.json}; Phase~2 runs under \texttt{nsys profile} and reuses the Inductor cache from Phase~1.

\subsection*{NVTX v3 Static Linking: ncu Range Mode Abandoned}

An appealing approach was to use \ncu's \texttt{--filter-by-nvtx-range} to replay only the kernels within each \texttt{aten::} NVTX range, collecting per-operator hardware counters without a full-workload replay.  PyTorch~2.x statically links NVTX~v3 inside \texttt{libtorch\_cuda.so}, whereas \ncu's injection hooks target the dynamically-loaded NVTX library.  As a result, \ncu's injection mechanism never observes any NVTX events; range filters silently capture nothing regardless of \texttt{emit\_nvtx} activity.  Application-mode replay (\texttt{--replay-mode application}) is the only path that reliably collects counters for Inductor-compiled workloads.

\subsection*{Functional vs.\ Aten IR Level Mismatch}

The initial backend implementation applied all FX passes at a single IR level.  Profiling the GPT-2 workload revealed that QKV fusion patterns were not found at runtime: after AOTAutograd decomposition, the shared \texttt{q/k/v} activation node had been replicated into per-consumer \texttt{aten.view} copies, destroying the three-GEMM pattern the fusion matcher expected.  This motivated the three-level architecture (\secref{sec:threelevel}).  The fix was non-trivial because the Dynamo functional level and the post-AOTAutograd Aten level look superficially similar; the difference only surfaces at profiling time when a supposedly-applied pass produces no measurable counter change.

\subsection*{FakeTensor Opacity in \texttt{\_aten\_inner\_compile}}

The \texttt{example\_inputs} argument of Inductor's \texttt{inner\_compile} hook may be \texttt{FakeTensor} objects---symbolic values that carry shape and dtype metadata but have no allocated storage.  An initial implementation of BatchNorm folding attempted to read \texttt{.data} from these tensors and crashed with storage-access errors.  The fix threads \texttt{real\_inputs = list(example\_inputs)} from the outer backend entry point before any \texttt{compile\_fx} call, so weight-reading passes operate on concrete tensors while Inductor's trace receives the \texttt{FakeTensor} inputs it expects.

\subsection*{GPU Clock Drift}

Initial example profiles were captured without clock management.  Baseline and optimized captures made at different times---with the optimized workload running cooler due to fewer scalar SIMT FLOPs---showed spurious duration differences on operators that were architecturally unchanged.  In one case the apparent speedup on an unmodified kernel exceeded 1.2$\times$, solely from clock drift.  The probe-and-lock mechanism (\secref{sec:clocklocking}) was added after this was identified as a systematic bias in the before/after comparison.

\subsection*{NVTX Warmup Attribution Contamination}

The layer deduplication NVTX wrapping initially emitted \texttt{layer::unique::} and \texttt{layer::duplicate::} ranges for every forward pass, including warmup iterations.  \texttt{ManifestBuilder}'s \texttt{\_tag\_layer\_partitions} attributed kernels launched during warmup to their enclosing layer-partition ranges, introducing warmup-only compilation artifacts into operator profiles and inflating attribution coverage.  The fix gates range emission behind an \texttt{\_NVTX\_STATE[`active']} flag set only after warmup completes.  Both bugs (clock drift and warmup contamination) were discovered concurrently and corrected in the same pipeline revision.

\subsection*{Kernel Invocation Count Invariant}

Invocation-order matching (\algref{alg:iom}) is correct only when the \texttt{--warmup-iters} and \texttt{--measure-iters} counts are identical between the \nsys capture and the \ncu replay.  A mismatch shifts kernel invocation indices, causing hardware counters to be assigned to the wrong operator call without any runtime error.  This invariant was violated during early development when a second \nsys capture used different iteration counts without regenerating the manifest; the resulting \texttt{profile.json} contained plausible-looking but incorrect counter values.  The pipeline now validates iteration counts before initiating replay.

\subsection*{sudo Environment Stripping Under \texttt{env\_reset}}

\ncu requires elevated privileges to access GPU performance counters.  Running \texttt{sudo ncu} on a system with the default \texttt{sudoers env\_reset} rule strips \texttt{PYTHONPATH} from the subprocess environment; the workload script imports from \texttt{nvidia.operator\_profiler}, so the subprocess exits with \texttt{ModuleNotFoundError} with no counter output and no obvious connection to the import path.  The fix uses \texttt{sudo env KEY=VAL \ldots\ ncu} to forward environment variables explicitly, bypassing the \texttt{env\_keep} list entirely.
